Teoremi e definizioni per l'orale:

Errore:
Teorema 2.2.1 su errore di rappresentazione e precisione di macchina (con dimostrazione)
Teorema 3.1.1. Errore totale somma di errore inerente e algoritmico (con dimostrazione)

Norme + localizzazione di autovalori:
Definizione 4.1.1 di norma vettoriale
Definizione 4.1.2 di norma matriciale
Definizione 4.1.3 di norma matriciale indotta da una norma vettoriale
Teorema 4.1.2 su compatibilità tra norma vettoriale e matriciale indotta 
Teorema 4.1.3 norme matriciali 1, 2, e infinito (solo enunciato)
Teorema sul condizionamento di un sistema lineare (non e' espresso come un teorema ma il risultato è a pagina 21 delle dispense Gemignani)
Teorema 4.4.1 di Gershgorin (con dimostrazione)
Secondo teorema di Gershgorin (solo enunciato, pag 81 dispense Bevilacqua-Menchi) 
Teorema 4.4.2 di Hirsh (con dimostrazione)

Sistemi lineari
Definizione 5.1.2 di matrice fattorizzabile LU
Teorema 5.1.1 di esistenza ed unicità della fattorizzazione LU (con dimostrazione)
Definizione 5.2.1 di matrici elementari di Gauss e proprietà a pgaina 31-32
Teorema 6.1.1 (fa la cosa giusta, con dimostrazione)
Definizione 6.1.1 di metodo convergente
Teorema 6.1.2 su condizioni sufficienti per la convergenza di un metodo iterativo per sistemi lineari (con dimostrazione)
Teorema 6.1.3 su condizioni necessarie per la convergenza di un metodo iterativo per sistemi lineari (con dimostrazione)
Teorema 6.1.5 su condizioni necessarie e sufficinenti per la convergenza di un metodo iterativo per sistemi lineari (con dimo solo nel caso P diagonalizzabile)
Definizione 6.3.1 di predomaninaza diagonale (per righe e per colonna)
Teorema 6.3.1 sulla convergenza di Jacobi e Gauss-Seidel sotto le ipotesi di predominanza diagonale (con dimo) + predominanza diag implica invertibilità (con dimostrazione)

Metodi per l'approssimazione di zeri di funzioni non lineari
Teorema 10.1.1 su metodo di bisezione (solo enunciato)
Teorema 10.2.1 (fa la cosa giusta, con dimo)
Definizione 10.2.1 su convergenza locale
Teorema 10.2.2 del punto fisso (con dimo)
Teorema 10.2.3 corollario al teorema del punto fisso
Definizione 10.2.2+3+4. di ordine di convergenza lineare, quadratica
Teorema 10.3.1 Convergenza locale del metodo delle tangenti e convergenza quadratico nel caso di radice semplice (con dimo)
Teorema 10.3.2 Convergenza in largo (aggiungere rispetto a quello scritto nelle dispense che sotto le ipotesi del teorema le successioni risultano monotone) senza dimo
 
Fattorizzazione QR: Matrici elementari di Householder (non per i compitini ma per l'orale degli appelli regolari : appunti caricati su elearning)
Minimi quadrati lineari e regressione lineare (non per i compitini ma per l'orale degli appelli regolari : appunti caricati su elearning)

